\begin{helper}
Let \(N\) be a natural number, and let
\(c,b\in\mathbb R^{\operatorname{Fin}(N)}\) have nonnegative integer
entries and equal total sums.  Suppose \(b\) is nonincreasing.  Suppose
there is a nonincreasing rearrangement \(x\) of \(c\) such that, for every
\(0\le r\le N\), the sum of the first \(r\) entries of \(x\) is at least
the sum of the first \(r\) entries of \(b\).  If
\(h\colon\mathbb R^{\operatorname{Fin}(N)}\to\mathbb R\) is convex and
symmetric under coordinate permutations, then
\[
h(b)\le h(c).
\]
\end{helper}
\begin{proof}
It is enough to prove that \(b\) belongs to the convex hull of the
permutations of \(x\).  Indeed, if
\[
b=\sum_{\sigma\in S_N} w_\sigma\,\sigma(x),\qquad
w_\sigma\ge0,\qquad \sum_{\sigma\in S_N}w_\sigma=1,
\]
then convexity gives
\[
h(b)\le \sum_{\sigma\in S_N}w_\sigma h(\sigma(x)).
\]
By symmetry each \(h(\sigma(x))\) is equal to \(h(x)\), so the right hand
side is \(h(x)\).  Since \(x\) is a permutation of \(c\), symmetry again
gives \(h(x)=h(c)\).  Thus the displayed convex-hull inclusion implies the
claim.

We now prove that convex-hull inclusion.  Write the prefix slack of a sorted
vector \(u\) relative to \(b\) as
\[
S_r(u)=\sum_{i<r}u_i-\sum_{i<r}b_i\qquad(0\le r\le N).
\]
The hypotheses say \(S_r(x)\ge0\) for all \(r\), and the equality of total
sums says \(S_N(x)=0\).  If \(x=b\), there is nothing to prove.  Otherwise,
let \(i\) be the first coordinate at which \(x_i\ne b_i\).  Since all
earlier coordinates agree, \(S_{i+1}(x)=x_i-b_i\).  The prefix inequality
forces \(x_i>b_i\); it cannot be \(x_i<b_i\).  Because the total discrepancy
is zero, some later coordinate must satisfy \(x_j<b_j\).  Let \(j>i\) be the
first such coordinate.

For every prefix ending between \(i\) and \(j-1\), a positive discrepancy at
\(i\) has appeared and no negative discrepancy has yet canceled it, so the
slack is strictly positive.  Define
\[
\delta=\min\Bigl(x_i-b_i,\ b_j-x_j,\ \min_{i<r\le j}S_r(x)\Bigr).
\]
All quantities in this minimum are positive integers: the entries of
\(b\) and \(x\) are nonnegative integers, and \(x\) is a permutation of
the nonnegative-integer vector \(c\).  Hence \(\delta\) is a positive
integer.

Let \(y\) be obtained from \(x\) by moving \(\delta\) units from coordinate
\(i\) to coordinate \(j\):
\[
y_i=x_i-\delta,\qquad y_j=x_j+\delta,\qquad y_k=x_k\ (k\ne i,j).
\]
The choices \(\delta\le x_i-b_i\) and \(\delta\le b_j-x_j\) show that the
affected coordinates remain nonnegative integers.  The total sum is
unchanged.  The prefix slacks are unchanged before \(i\), reduced by
\(\delta\) for the prefixes that contain \(i\) but not \(j\), and unchanged
again after \(j\).  Since \(\delta\) was chosen no larger than all the
intermediate slacks, every prefix slack of \(y\) is still nonnegative.

This transfer also has the needed convex-combination form.  Because both
\(x\) and \(b\) are nonincreasing and \(i<j\), we have \(b_i\ge b_j\); from
\(x_i>b_i\) and \(x_j<b_j\) it follows that \(x_i>x_j\).  Put
\[
\theta=\frac{\delta}{x_i-x_j}.
\]
The denominator is positive, and
\[
x_i-x_j=(x_i-b_i)+(b_i-b_j)+(b_j-x_j)
\]
with all three summands nonnegative and the first and third positive.  Since
\(\delta\le x_i-b_i\) and \(\delta\le b_j-x_j\), we get
\(0\le\theta\le1\).  On the two affected coordinates,
\[
(y_i,y_j)
=(1-\theta)(x_i,x_j)+\theta(x_j,x_i),
\]
and every other coordinate is unchanged.  Thus
\[
y=(1-\theta)x+\theta\,\tau(x),
\]
where \(\tau\) is the transposition swapping \(i\) and \(j\).

The vector \(y\) may fail to be nonincreasing.  Let \(y^\downarrow\) be a
nonincreasing rearrangement of \(y\).  For each \(r\), the sum of the first
\(r\) entries of \(y^\downarrow\) is the maximum possible sum of \(r\)
coordinates of \(y\), so it is at least the sum of the first \(r\)
coordinates in the displayed ordering of \(y\).  Hence
\(y^\downarrow\) still has all prefix slacks nonnegative.  Since
\(y^\downarrow\) is a permutation of \(y\), applying the sorting permutation
to the convex-combination identity above shows
\[
y^\downarrow
=(1-\theta)\rho(x)+\theta\,\rho(\tau(x))
\]
for some permutation \(\rho\).  Therefore \(y^\downarrow\) lies in the
convex hull of the permutations of \(x\).

It remains to see that iterating this operation reaches \(b\).  Use the
integer measure
\[
D(u)=\sum_{k=1}^N |u_k-b_k|
\]
on nonincreasing nonnegative-integer vectors \(u\) with the same total sum
as \(b\) and with nonnegative prefix slacks.  Before resorting, the transfer
decreases the two affected absolute values by \(\delta\) each:
\[
|y_i-b_i|=|x_i-b_i|-\delta,\qquad
|y_j-b_j|=|x_j-b_j|-\delta,
\]
and all other terms are unchanged.  Thus \(D(y)=D(x)-2\delta<D(x)\).
Sorting cannot increase \(D\) relative to the already sorted vector \(b\).
Indeed, an adjacent inverted pair in \(y\) has values \(a\le d\) in the
order \((a,d)\), while the corresponding entries of \(b\) satisfy
\(b_r\ge b_s\).  The elementary rearrangement inequality
\[
|d-b_r|+|a-b_s|\le |a-b_r|+|d-b_s|
\]
shows that swapping the larger value leftward does not increase the
two-coordinate contribution to \(D\).  Repeating adjacent swaps until the
vector is sorted gives
\[
D(y^\downarrow)\le D(y)<D(x).
\]

Since \(D\) is a nonnegative integer, this strictly decreasing process
terminates.  A terminal vector satisfying the same hypotheses must be
\(b\): if it were not \(b\), the first positive discrepancy \(i\), the first
later negative discrepancy \(j\), and the positive transfer just constructed
would still be available.  Therefore, starting from \(x\), finitely many
steps reach \(b\), and at each step the new sorted vector lies in the convex
hull of permutations of the previous one.  Composing these finite convex
combinations shows that \(b\) lies in the convex hull of the permutations of
\(x\), as required.
\end{proof}
